% 本文件由 03_代码/p2_lambda_tables.py 自动生成，请勿手工编辑。
% 数据来源：06_支撑材料/p2_lambda_ablation.json, p2_lambda_diag.json,
%           p2_results.json

\section{问题二日末储备惩罚的消融与库存口径检验}
\label{app:p2-lambda}

日前目标函数中的 $\lambda\xi$ 项定价的是\textbf{计划轨迹} $\bar E$ 的日末值缺口，而不是实际执行后的储电量。标定在 $1$ 月选中 $\lambda=0.00$并冻结整年，此时 $\xi$ 在经济上完全自由，终端约束不再产生定价信号。本节把这件事实测清楚，并检验它是否只是评价口径造成的假象。

\subsection{参考轨迹的退化程度}

按交付口径（$2$ 月 $1$ 日--$12$ 月 $31$ 日共 $334$ 天，$\rho=0.9$）统计，参考轨迹 $\bar E$ 贴上限 $E_{\max}=10800$~kWh 的时段占 $15.64\%$、贴下限 $E_{\min}=1200$~kWh 的时段占 $18.23\%$，而\textbf{日末值贴下限的天数占比为 $100.00\%$}，日末均值恰为 $1200.00$~kWh。也就是说 $\bar E$ 在日内仍是一条有形态的轨迹（均值 $6,243.72$、标准差 $3,709.35$~kWh），但\textbf{其日末值已不含任何跨日信息}，$\xi_n\equiv E^{\mathrm{tar}}-E_{\min}=4800$~kWh 逐日恒定。

这不影响执行层的合法性：实际储电量 $E$ 仍由因果规则给出、逐段满足物理约束，其均值 $6,845.83$~kWh、贴下限时段占比仅 $1.20\%$。储备线 $R_t=E_{\min}+\rho(\bar E_t-E_{\min})$ 是\textbf{计划跟随}约束而非跨日储能策略：$\rho=0.9$ 时实际电量低于 $R_t$ 的时段占 $9.13\%$，其中 $7.81\%$ 因储备线而未能放电（含 $6.72\%$ 的时段储电量仍在 $3000$~kWh 以上），而 $\rho=0.5$ 的对照臂该占比仅 $1.01\%$，说明 $R_t$ 确实在起作用；只是在 $\bar E$ 日末值退到下限后，$R_t$ 对尾部时段的约束力随之减弱。

\subsection{只改 $\lambda$ 的全年消融}

\begin{table}[htbp]
  \centering
  \caption{日末储备惩罚 $\lambda$ 的全年消融：固定 $(\alpha,\rho)$、关闭滚动重标定，只改 $\lambda$（2025-02-01--12-31，共 334 天）}
  \label{tab:p2-lambda}
  \footnotesize
  \begin{tabular}{ccccrrrr}
    \toprule
    $\alpha$ & $\rho$ & $\lambda$ & 费用/万元 & 期末储电量/kWh & 紧急购电量/MWh & $\bar E$ 日末/kWh & $E<R_t$ 占比 \\
    \midrule
    0.75 & 0.9 & 0.00 & \textbf{1,389.13} & \textbf{1,559.7} & 175.08 & \textbf{1,200.0} & 7.65\% \\
    0.75 & 0.9 & 0.20 & \textbf{1,389.13} & \textbf{1,559.7} & 175.08 & \textbf{1,200.0} & 7.65\% \\
    0.75 & 0.9 & 0.50 & 1,389.27 & 6,174.3 & 164.58 & 6,000.0 & 6.62\% \\
    0.75 & 0.9 & 0.80 & 1,389.27 & 6,174.3 & 164.58 & 6,000.0 & 6.62\% \\
    0.75 & 0.9 & 1.20 & 1,389.27 & 6,174.3 & 164.58 & 6,000.0 & 6.62\% \\
    \midrule
    0.75 & 0.5 & 0.00 & \textbf{1,394.75} & \textbf{1,559.7} & 165.74 & \textbf{1,200.0} & 1.01\% \\
    0.75 & 0.5 & 0.20 & \textbf{1,394.75} & \textbf{1,559.7} & 165.74 & \textbf{1,200.0} & 1.01\% \\
    0.75 & 0.5 & 0.50 & 1,394.43 & 6,174.3 & 151.07 & 6,000.0 & 0.96\% \\
    0.75 & 0.5 & 0.80 & 1,394.43 & 6,174.3 & 151.07 & 6,000.0 & 0.96\% \\
    0.75 & 0.5 & 1.20 & 1,394.43 & 6,174.3 & 151.07 & 6,000.0 & 0.96\% \\
    \midrule
    0.70 & 0.9 & 0.00 & \textbf{1,395.36} & \textbf{1,421.7} & 225.34 & \textbf{1,200.0} & 8.93\% \\
    0.70 & 0.9 & 0.20 & \textbf{1,395.36} & \textbf{1,421.7} & 225.34 & \textbf{1,200.0} & 8.93\% \\
    0.70 & 0.9 & 0.50 & 1,394.43 & 6,059.1 & 212.17 & 6,000.0 & 8.13\% \\
    0.70 & 0.9 & 0.80 & 1,394.43 & 6,059.1 & 212.17 & 6,000.0 & 8.13\% \\
    0.70 & 0.9 & 1.20 & 1,394.43 & 6,059.1 & 212.17 & 6,000.0 & 8.13\% \\
    \bottomrule
  \end{tabular}
  \par\vspace{2pt}
  \begin{minipage}{.94\textwidth}\footnotesize
    注：\textbf{加粗}表示该臂的参考轨迹 $\bar E$ 日末值在全部 $334$ 天都贴在下限 $E_{\min}=1200$~kWh，即 $\xi$ 项完全失效。三组 $(\alpha,\rho)$一致地显示：$\lambda\le0.20$ 时 $\bar E$ 日末恒为 $1200$~kWh，$\lambda\ge0.50$ 时恒为 $6000$~kWh，且在 $\lambda\ge0.5$ 的三个取值上费用差异不超过 $2$~元——终端目标已从软惩罚实际变成硬约束。同一组内 $\lambda$ 由 $0$ 提到 $0.50$，报告区间费用变化为 $-0.067\%\sim+0.010\%$，属同一量级的数值噪声。
  \end{minipage}
\end{table}

消融同时给出期末储电量：$\lambda\le0.20$ 时全年结束时储能被消耗到 $1,421.7$--$1,559.7$~kWh，$\lambda\ge0.50$ 时保留 $6,059.1$--$6,174.3$~kWh（区间端点是 $(\alpha,\rho)$ 的三个取值之间的差异）。这说明 $\lambda=0$ 确实带来把库存用得更空的效果，但把这一库存差异按元/kWh 折价后（$v\in[0.35,0.65]$），三组 $(\alpha,\rho)$ 的费用排序均不变。

\subsection{一月选参是否被库存口径扭曲}

$1$ 月选参的准则是窗口内总费用最小，而窗口截断会系统性偏爱\textbf{消耗库存}的策略：少买的电计入费用，少掉的库存却不计价。把准则换成库存修正费用 $J+v\,(E_{\text{初}}-E_{\text{末}})$ 后，$\lambda=0$ 还是不是最优，取决于库存单价 $v$。下表在 $1$ 月选参的全部候选点上扫描 $v$。

\begin{table}[htbp]
  \centering
  \caption{库存计价准则 $J+v\,(E_{\text{初}}-E_{\text{末}})$ 下一月选参最优解对库存单价 $v$ 的敏感性}
  \label{tab:p2-breakeven}
  \footnotesize
  \begin{tabular}{cccc}
    \toprule
    库存单价 $v$/(元$\cdot$kWh$^{-1}$) & 最优 $\lambda$ & 最优 $(\alpha,\rho)$ & 库存修正后一月费用/元 \\
    \midrule
    0.00 & 0.00 & $(0.75,\ 0.9)$ & 2,219,466.61 \\
    0.35 & 0.00 & $(0.75,\ 0.9)$ & 2,220,765.92 \\
    0.50 & 0.00 & $(0.75,\ 0.9)$ & 2,221,322.77 \\
    0.65 & 0.00 & $(0.75,\ 0.9)$ & 2,221,879.61 \\
    1.00 & 0.00 & $(0.75,\ 0.9)$ & 2,223,178.93 \\
    1.50 & 0.00 & $(0.75,\ 0.9)$ & 2,225,035.09 \\
    2.00 & \textbf{0.50} & $(0.75,\ 0.9)$ & 2,225,572.74 \\
    \midrule
    \multicolumn{4}{l}{\footnotesize 在 $[0,4]$~元/kWh 上二分求得翻转点 $v^*=1.712$~元/kWh} \\
    \bottomrule
  \end{tabular}
  \par\vspace{2pt}
  \begin{minipage}{.94\textwidth}\footnotesize
    注：$E_{\text{初}}=6000$~kWh，$E_{\text{末}}$ 取 $2$ 月 $1$ 日储电量，候选点共 $180$ 个。$v$ 的含义是期末少存 $1$~kWh 相当于省下多少钱。储电量在最贵时段顶替计划购电的边际价值上界为 $\eta_c\eta_d\max_t p_t=1.130$~元/kWh（电价区间 $[0.3713,$~$1.3952]$~元/kWh）；$v^*/$上界 $=1.51$，即只有在把库存单价高估到其可实现边际价值的 $1.5$ 倍以上时，标定才会改选 $\lambda>0$。故 $\lambda=0$ \textbf{是稳健的标定结果}，而不是评价口径的假象。
  \end{minipage}
\end{table}

\subsection{结论}

$\lambda=0.00$ 是\textbf{标定结果而非建模缺陷}：它使参考轨迹 $\bar E$ 的日末值退化为下限，这是该取值下的必然表现，已在正文如实说明；但（i）把它提高到 $0.5$ 以上对全年费用的影响在 $-0.067\%\sim+0.010\%$ 之间，（ii）在库存计价准则下，只有把库存单价高估到其可实现边际价值的 $1.5$ 倍以上，标定才会改选 $\lambda>0$。因此\textbf{$\xi$ 项的存在使是否保留日末储备成为一个由数据回答的问题}，而不是一个未加检验的隐含假设——这正是引入该项的意义所在。
